% Appendix F: drafting and repair hypotheses.
% Checkpoint facts: config and weights at the pinned revisions. SGLang's DFlash cycle: code
% reading at the pin. Byte counts: derived. Card numbers (B200): the authors' claims.
% Acceptance by position and domain (correctness runs, no timings): register entries
% drafter-acc, drafter-trace, drafter-launch. Cycle timing, equality, training, P6, P2, P3:
% hypotheses or designs; every measurement pending. Proposition prefix-value and the
% finite-map proposition: proved (no novelty claimed).
\section{Drafting and repair hypotheses}\label{app:drafting}\label{sec:drafting}
How much a verify pass is worth depends on the drafter in front of it: a cycle commits as many tokens as the target accepts plus one, and its cost includes the draft. The native MTP layer of Qwen3.5-4B drafts sequentially, paying one small layer plus the full 1.27~GB head per drafted token, and its acceptance decays with depth. A block drafter proposes a whole block in one forward pass. Every drafter that projects through a head pays for it at every draft step: Medusa and the EAGLE family~\cite{medusa,eagle,eagle2,eagle3}, native multi-token prediction~\cite{gloeckle2024,deepseekv3} and DFlash~\cite{dflashpaper}. This appendix describes the public block drafter for this target, how SGLang runs it on a hybrid model and what is measured about it, and then states three proposals as hypotheses with their kill tests: a selector trained for rate (P6) and exact repair of long draft windows (P2, P3). None of the proposals has a result yet.

\subsection{The public block drafter}\label{app:drafter-model}\label{sec:drafter-model}
DFlash has a public drafter for exactly this target, \code{z-lab/Qwen3.5-4B-DFlash} at revision \texttt{9a1996cc}, under Apache-2.0~\cite{dflashpaper,dflash4bcard}; \code{modal-labs/Qwen3.5-4B-DFlash} carries byte-identical weights. Read from the checkpoint rather than the card, it is a six-layer Qwen3-style transformer of width 2,560 (32 query and 8 key--value heads of dimension 128, MLP width 9,216) with no embedding or head of its own: it embeds with the target's table and projects through the target's tied head. It conditions on the target's hidden states after layers 1, 5, 9, 13, 17, 21, 25 and 29, concatenated per token and projected by one linear map and an RMS norm to a single context vector. Five layers use sliding-window attention (window 4,096) and the last uses full attention. It has 634 million parameters, 1.27~GB in BF16, almost exactly the size of the target's head (derived).

A cycle in SGLang's DFlash worker~\cite{sgcode} works as follows (code reading at the pin). The block is the last committed token followed by fifteen copies of a mask token (id 248,077). Each draft layer's keys and values for the committed context come from the projected target features, which the worker writes into the drafter's own cache as tokens are committed. One forward pass over the sixteen positions yields sixteen hidden states; positions 1--15 go through the target head, and their argmaxes are the drafted tokens. The target verifies the block in one pass, the worker accepts the longest drafted prefix that matches the target's argmax, commits it plus the target's token at the first mismatch, and appends the committed tokens' features to the drafter's cache. The sliding layers are causal inside the block and the full layer is not; SGLang infers this from the layer types because the configuration sets no \code{is\_causal}, which matches the attention pattern of the DFlash training recipe in SGLang's SpecForge repository (\code{sgl-project/SpecForge}).

On the hybrid target the verify pass has a cost the drafter's parameters do not show. SGLang verifies a chain on the GDN layers by keeping the recurrent state after every block position, so that the state after the accepted prefix can be committed without recomputation. Each state is 50.3~MB per request, so a 16-token block writes about 805~MB per request per cycle, and each running request reserves space for sixteen of them (derived). That bounds block-16 capacity on this machine, and it is a property of the verifier, not of the drafter.

The card's launch line targets a B200 with TRT-LLM target attention, FA4 draft attention and piecewise-compiled prefill graphs. None carries over unchanged: TRT-LLM attention is Blackwell-only, FA3 and FA4 are unavailable on this aarch64 host, and the piecewise prefill graph fails to capture for this target at the pinned commit. The drafter is served with FlashInfer for target and draft attention, the FlashInfer GDN kernels, the default prefill graph and the card's overlap planning stream, with decode, verify and draft CUDA graphs and the overlap scheduler on; the launch records of the acceptance runs below give the exact command, environment and resolved server settings~\ev{drafter-launch}.

\subsection{What a cycle costs}\label{app:drafter-cost}\label{sec:drafter-cost}
Counting weight bytes gives the shape of the trade before any timing. At batch one a DFlash cycle reads the drafter once (1.27~GB), the target head once for the fifteen drafted positions (1.27~GB), and the target once to verify (8.41~GB): drafting costs about 30\% of the verify pass's weight traffic and proposes fifteen tokens. An MTP chain of three steps reads the MTP layer and the full head three times (about 4.5~GB, 54\% of the verify pass) and proposes three (derived). Activation traffic, the per-position GDN states, attention over the context and host time are not in these counts; the measured MTP cycle of Appendix~\ref{app:profile} shows how much the host adds. The split of the DFlash cycle at batch 1, 4, 16 and 64 into draft forward, draft head, verification, GDN commit, cache update and host time is \pend{drafter-cycle}.

\subsection{Acceptance by block position}\label{app:drafter-acceptance}\label{sec:drafter-acceptance}
Let $h_j$ be the number of verify cycles in which exactly $j$ drafted tokens were accepted, $N=\sum_jh_j$, and $S(k)=\sum_{j\ge k}h_j/N$ the fraction of cycles that accept at least $k$. The conditional acceptance of block position $k$ is $\alpha_k=S(k)/S(k-1)$, and a cycle commits $\tau=1+\sum_{k\ge1}S(k)$ tokens on average over cycles. SGLang reports $h$ per request, so $\alpha_k$ and $\tau$ are measured rather than modelled. The model card instead averages each request's tokens per cycle over requests; we call that quantity $\bar\tau_{\rm req}$. It weights short responses more than $\tau$ does, and for these runs it is the larger of the two. The card reports mean accepted lengths of 5.9--7.7 at block 16 on a B200 for its five benchmarks, against 3.3--3.5 for MTP with three steps~\cite{dflash4bcard}; The same card reports, on one B200 with greedy decoding, 3.40--4.60$\times$ over autoregressive decoding at concurrency~1 with blocks of 16 and 2.15--2.61$\times$ at concurrency~32 with blocks of 8, against at most 2.31$\times$ and 1.80$\times$ for the MTP layer. These are the authors' numbers on other hardware.

The measurement uses a panel of 32 MATH-500 problems and the first 16 chat, 16 code and 16 GSM8K prompts of the benchmark's earlier confirmation split, decoded greedily with thinking on, stopping naturally at up to 2,048 new tokens, at concurrency 1, with blocks of 16 and 8~\ev{drafter-acc}. The runs used an engine trace hook that synchronizes every cycle on a shared GPU, so acceptance and token counts are valid and timings are not; the panel contains benchmark test problems and measures acceptance only, never quality.

\begin{table}[htbp]
\small\centering
\begin{tabular}{@{}llrrrrr@{}}\toprule
Block & Average & All & Chat & Code & GSM8K & MATH-500\\\midrule
16 & $\tau$, over cycles & 6.18 & 3.98 & 7.63 & 6.72 & 7.55\\
16 & $\bar\tau_{\rm req}$, over requests (the card's) & 6.98 & 4.26 & 7.87 & 7.16 & 7.81\\
8 & $\tau$, over cycles & 5.04 & 3.64 & 5.75 & 5.31 & 5.77\\
8 & $\bar\tau_{\rm req}$, over requests & 5.34 & 3.80 & 5.76 & 5.45 & 5.84\\\midrule
\multicolumn{2}{@{}l}{Requests (verify cycles at block 16)} & 80 (21,274) & 16 & 16 & 16 & 32\\\bottomrule
\end{tabular}
\caption{Committed tokens per verify cycle for the public DFlash drafter on GH200 (greedy, thinking on, concurrency 1), pooled over cycles ($\tau$) and averaged over requests as the model card does ($\bar\tau_{\rm req}$). Measured in correctness runs~\ev{drafter-acc}; the survival form $1+\sum_kS(k)$ agrees with the pooled $\tau$ to within 0.01.}
\label{tab:drafter-tau}
\end{table}

A block of 16 commits 6.2 tokens per verify cycle and a block of 8 commits 5.0 (Table~\ref{tab:drafter-tau}). Acceptance depends strongly on the domain: chat commits about four tokens per cycle, GSM8K 6.7, and code and MATH-500 about 7.6. The per-request MATH-500 average, 7.8, sits near the 7.48 the card reports for that benchmark in the same convention~\cite{dflash4bcard}, but that is a cross-hardware reference, not a reproduction: the card used a B200 and a 4,096-token cap and does not state how many problems it used, where this panel uses a GH200, a 2,048-token cap and 32 sampled problems.

The per-position rates of Figure~\ref{fig:drafter-acceptance} show why a longer block still commits more tokens per cycle; whether it pays depends on the cost of a cycle, which is pending. Over all domains the conditional acceptance dips at the second position (0.80) and then rises toward 0.9, so the share of cycles that reach position 15 is still 13\%. The rise is a selection effect: cycles that get that far are in predictable stretches of text, which is not evidence that the drafter is better deep in the block. Chat is lower throughout (0.69 at positions 2 and 3), and the dip at position 2 appears in every domain. Two further observations bound what training could change. Acceptance, pooled over tokens, is lower in the answer after the thinking text than in the thinking itself for chat (3.3 against 4.2 tokens per cycle) and higher for code and mathematics, so no single segment dominates. And the first rejected draft in a cycle equals the target's token at the previous position in about 6\% of cycles, a repetition-type failure, while in 94\% it is some other token~\ev{drafter-trace}. The committed tokens per cycle against tuned MTP on the same prompts are \pend{drafter-cycle}.

\begin{figure}[htbp]
\centering
\datafigure{../evidence/drafter/acceptance_by_position.csv}{%
\begin{tikzpicture}
\begin{groupplot}[group style={group size=2 by 1,horizontal sep=1.8cm},width=.45\linewidth,height=5.0cm,
  xmin=0.5,xmax=15.5,xtick={1,3,5,7,9,11,13,15},xlabel={Block position $k$},grid=major,unbounded coords=discard]
\nextgroupplot[ymin=0.6,ymax=1,ylabel={Conditional acceptance $\alpha_k$},title={(a) $\alpha_k=S(k)/S(k-1)$},
  legend to name=acclegend,legend columns=5,legend style={/tikz/every even column/.append style={column sep=6pt}}]
\addplot[thick,ink,mark=*,mark size=1.3pt,discard if not={domain}{all},discard if not={block_size}{16}] table[x=position,y=accept_prob,col sep=comma]{../evidence/drafter/acceptance_by_position.csv};\addlegendentry{all}
\addplot[semithick,ochre,dashed,mark=square,mark options={solid},mark size=1.4pt,discard if not={domain}{chat},discard if not={block_size}{16}] table[x=position,y=accept_prob,col sep=comma]{../evidence/drafter/acceptance_by_position.csv};\addlegendentry{chat}
\addplot[semithick,teal,mark=triangle,mark size=1.6pt,discard if not={domain}{code},discard if not={block_size}{16}] table[x=position,y=accept_prob,col sep=comma]{../evidence/drafter/acceptance_by_position.csv};\addlegendentry{code}
\addplot[semithick,brick,dashdotted,mark=diamond,mark options={solid},mark size=1.6pt,discard if not={domain}{math},discard if not={block_size}{16}] table[x=position,y=accept_prob,col sep=comma]{../evidence/drafter/acceptance_by_position.csv};\addlegendentry{GSM8K}
\addplot[semithick,muted,densely dotted,mark=o,mark options={solid},mark size=1.3pt,discard if not={domain}{math500},discard if not={block_size}{16}] table[x=position,y=accept_prob,col sep=comma]{../evidence/drafter/acceptance_by_position.csv};\addlegendentry{MATH-500}
\nextgroupplot[ymode=log,ymin=0.02,ymax=1,ylabel={Share of cycles reaching $k$, $S(k)$},title={(b) Survival $S(k)$},
  ytick={0.02,0.05,0.1,0.2,0.5,1},yticklabels={0.02,0.05,0.1,0.2,0.5,1}]
\addplot[thick,ink,mark=*,mark size=1.3pt,discard if not={domain}{all},discard if not={block_size}{16}] table[x=position,y=survival,col sep=comma]{../evidence/drafter/acceptance_by_position.csv};
\addplot[semithick,ochre,dashed,mark=square,mark options={solid},mark size=1.4pt,discard if not={domain}{chat},discard if not={block_size}{16}] table[x=position,y=survival,col sep=comma]{../evidence/drafter/acceptance_by_position.csv};
\addplot[semithick,teal,mark=triangle,mark size=1.6pt,discard if not={domain}{code},discard if not={block_size}{16}] table[x=position,y=survival,col sep=comma]{../evidence/drafter/acceptance_by_position.csv};
\addplot[semithick,brick,dashdotted,mark=diamond,mark options={solid},mark size=1.6pt,discard if not={domain}{math},discard if not={block_size}{16}] table[x=position,y=survival,col sep=comma]{../evidence/drafter/acceptance_by_position.csv};
\addplot[semithick,muted,densely dotted,mark=o,mark options={solid},mark size=1.3pt,discard if not={domain}{math500},discard if not={block_size}{16}] table[x=position,y=survival,col sep=comma]{../evidence/drafter/acceptance_by_position.csv};
\end{groupplot}
\node[anchor=north] at ($(group c1r1.south)!.5!(group c2r1.south)+(0,-1.0cm)$) {\pgfplotslegendfromname{acclegend}};
\end{tikzpicture}}{Conditional acceptance and survival by block position, block size 16, per domain.}
\caption{The public DFlash drafter by block position, block size 16, greedy with thinking on, per domain of the panel. (a)~Conditional acceptance $\alpha_k$; (b)~the share of verify cycles that accept at least $k$ drafted tokens, on a logarithmic axis. Measured in correctness runs~\ev{drafter-acc}; the late rise in (a) is a selection effect of the few cycles that get that far.}
\label{fig:drafter-acceptance}
\end{figure}

\subsection{Output equality and limits}\label{app:drafter-limits}\label{sec:drafter-limits}
Greedy speculative verification is exact in real arithmetic, but the verify pass computes the target's logits at a different batch shape from plain decoding, so the stock noise floor of Appendix~\ref{app:divergence} applies; DFlash's divergences against plain decoding per 1,000 tokens (blocks 8 and 16), classified by mechanism and set against the plain-decoding floor, are \pend{drafter-equality}. Four limitations are candidates: acceptance that decays with block position or by domain; drafting cost dominated by the full-vocabulary head for fifteen positions; the verifier's per-position state writes, which bound block size at high concurrency; and draft states far from the target's in the head's metric, which is what defeats transport on DFlash drafts (median relative drift 0.92, Section~\ref{sec:transport-measured}). The public drafter's training data are not published, so none of these limitations can be traced to its data. A custom drafter is worth training only against one of them, judged by total serving time: fine-tuning on the target's own outputs for acceptance, adding DFlash~2's selector~\cite{dflash,dflashcard} for acceptance, or an auxiliary loss on $\norm{h^t-h^d}_2^2/\norm{h^t}_2^2$ if the roughly hundredfold drift reduction transport needs proves reachable (Appendix~\ref{app:transport-scaling}). A training set of 18,000 prompts, disjoint from every benchmark and quality set by construction, is built and its manifest committed~\ev{drafter-data}; which limitation the measurements support, and the decision on training, are \pend{drafter-train}.

\subsection{Training a selector for rate (P6, hypothesis)}\label{app:selector}\label{sec:drafter-selector}
A later proposal keeps the DFlash backbone and the verifier frozen and trains only a small predecessor-conditioned selector over the backbone's $K=16$ candidates per position, $q_i(b\mid a)=\softmax_{b\in C_i}[\ell_i(b)+s_{\theta,i}(a,b)]$, so that one coherent path is proposed; DFlash~2's selector architecture~\cite{dflash}, whose training objective the announcement does not state, is the baseline, not a novelty. The objective is emitted tokens per unit time, $\E[\sum R]/\E[\sum C]$ with $R=L+1$ per cycle, trained as $R-\lambda C$ with $\lambda$ set to the measured rate. If $\lambda^\ast$ is the best achievable ratio in Equation~\eqref{eq:rate}, then $\E R-\lambda^\ast\E C\le0$ for every policy, with equality exactly for the policies that attain the ratio, because $\E C>0$; so maximizing the difference at the right $\lambda$ maximizes the rate. Two identities make the accepted-length term trainable exactly.
\begin{proposition}[Prefix value of a sampled draft policy]\label{prop:prefix-value}
Let the draft path $Y$ be sampled position by position, $Y_i\sim q_i(\cdot\mid Y_{<i})$, and let $L$ be the number of accepted drafted tokens. Against a greedy target with continuation $g$, $\E L=\sum_k\prod_{i\le k}q_i(g_i\mid g_{<i})$, and $\nabla_\theta\E L=\sum_j\bigl(\sum_{k\ge j}S_k\bigr)\nabla_\theta\log q_j(g_j\mid g_{<j})$ with $S_k=\prod_{i\le k}q_i(g_i\mid g_{<i})$. Against a sampled target with the standard rule, with $A_k=\prod_{i\le k}\min(1,p_i(Y_i)/q_i(Y_i))$,
\[\E L=\E_Y\Bigl[\sum_kA_k\Bigr],\qquad\nabla_\theta\E L=\E_Y\Bigl[\sum_j\ind\{q_j(Y_j)<p_j(Y_j)\}\Bigl(\sum_{k\ge j}A_k\Bigr)\nabla_\theta\log q_j(Y_j)\Bigr].\]
\end{proposition}
\begin{proof}
Greedy: $L\ge k$ exactly when $Y_i=g_i$ for $i\le k$, which has probability $S_k$, and $\E L=\sum_k\Pr(L\ge k)$; differentiate the product. Sampled: given the path, position $k$ is reached and accepted with probability $A_k$, so $\E L=\E_Y\sum_kA_k$. Its gradient has a score part and a pathwise part. In the score part, $\E[A_k\nabla\log q_j(Y_j)]$ vanishes for $k<j$, because $A_k$ depends only on $Y_{<j}$ and the score has conditional mean zero, leaving $\sum_j\nabla\log q_j\sum_{k\ge j}A_k$. In the pathwise part, $\nabla\min(1,p_j/q_j)$ is zero where $q_j<p_j$ and $-\min(1,p_j/q_j)\nabla\log q_j$ where $q_j\ge p_j$, which contributes $-\ind\{q_j\ge p_j\}\sum_{k\ge j}A_k\nabla\log q_j$. The two parts add to the stated expression.
\end{proof}
Freezing the acceptance probabilities as a fixed reward keeps only the score part and so gives a biased gradient. The proposition is a direct calculation and we claim no novelty for it. Acceptance-aware training already optimizes the greedy expected accepted length directly, with a window total-variation loss for sampling~\cite{aadt}; BV Loss trains on the exact expected acceptance length under block verification~\cite{bvloss}; D-PACE, PARD-2 and VAT weight a cross-entropy by position or by the chance of reaching it~\cite{dpace,pard2,vat}; variational speculative decoding optimizes a bound on sequence acceptance~\cite{vsd}; TIGER uses the realized accepted length as a reinforcement-learning reward~\cite{tiger}; and DBLAST drafts blocks with explicit dependence between positions~\cite{dblast}. The only trained DFlash selector we know of, LiLiCorr's, uses cross-entropy over candidates~\cite{lilicorr}. What remains for P6 is the combination: a DFlash~2-style selector over a frozen backbone, trained for rate rather than for accepted length, with the sampled-target gradient above; whether that pays is the experiment. The verification width must be chosen from committed-prefix features before proposals are drawn, for the reason of Appendix~\ref{app:stopping}; choosing a draft-tree budget from a profiled verify-cost curve at inference, as CaDDTree does for DFlash marginals~\cite{caddtree}, and confidence-scheduled width, as in DSpark~\cite{dspark}, are the existing ways to set it.

The cheapest test needs no training. From ordinary greedy runs, log the candidate sets $C_i$ and the completed greedy continuation $g$, and compute $U=\max\{k:g_i\in C_i\text{ for all }i\le k\}$, an upper bound on the accepted length of any selector over those frozen candidates; if $U$ barely exceeds DFlash's actual acceptance there is little to gain. Only then is the rate objective compared, at fixed width, with the strongest matched objective (VAT's weighting ranks first in its own ablation~\cite{vat}; D-PACE-style weighting and position-weighted cross-entropy are the other controls) using the same backbone, candidates, selector size, features, data and budget. The pre-registered effect size is $1.25\times$ end to end over the strongest matched baseline, rejected if the upper confidence bound falls below it; a gain that vanishes against an identical selector trained with another objective supports the selector, not the objective. The candidate-support screen $U$ against measured acceptance, and then the matched-objective comparison, are \pend{drafter-p6}.

\subsection{Exact repair of long windows (P2, P3, hypotheses)}\label{app:repair}\label{sec:repair}
A block drafter proposes a whole block in one forward pass and the target verifies it in one pass; the public DFlash card reports a mean accepted length of 7.48 tokens per verify pass at block size 16 on MATH-500 on a B200~\cite{dflash4bcard}. While decoding is limited by weight traffic, a verify pass over a wider block costs little more than one over a narrow block, so the question is whether the target can commit far more than about seven tokens per pass without giving up exactness. Two proposals are tested with oracle experiments designed to kill them cheaply: long-window repair of a DFlash window by the target's own predictions (P2), and target-anchored residual decoding (P3), which repairs a window by updating the target's previous computation with fixed low-rank corrections and then audits it exactly. Both keep exactness by the same device: only tokens that an exact audit pass of the target has verified are committed, approximate states never enter the committed cache, and the noise of a position is never resampled.

\paragraph{Prior work} Reusing the verifier's work on rejected positions is an active line. ReTrace conditions the next DFlash block on the rejected trajectory~\cite{retrace2026}, DFlow feeds the verifier's hidden states at rejected positions back to the drafter~\cite{dflow}, Carryover Drafting reuses those states as extra draft context~\cite{carryover2026}, and CURE repairs uncertain positions of block-parallel drafts locally~\cite{cure2026}; token recycling reuses candidate tokens from earlier steps~\cite{tokenrecycling2024}, and CARD and H-Spec restructure parallel speculation around cached queries and drafter state~\cite{card2025,hspec2026}. Jacobi-style iteration has been applied to translation, to feedforward computation, to image generation with probabilistic acceptance and to causal language models trained for it~\cite{santilli2023jacobi,song2021parallel,sjd2025,so2026scd,hu2026jacobiforcing}, and speculative speculative decoding overlaps drafting with verification~\cite{specspec}. P3's anchored update is closest to computation reuse across similar inputs, which DeltaCNN, Eventful Transformers and CacheBlend apply to video frames, vision tokens and cached retrieval contexts~\cite{deltacnn2022,eventful2023,cacheblend2025}. None of these, as far as we found, commits a long window through an exact audit of a cheap residual evaluator of the target itself; that combination is what the kill tests below examine.

\paragraph{Setting and cost model}\label{sec:repair-model} Greedy decoding solves a triangular system: with $F_i(y_{<i})$ the target's argmax after the committed prefix and $y_1,\dots,y_{i-1}$, the continuation satisfies $y_i=F_i(y_{<i})$. A verify pass evaluates $F$ at every position of a candidate block from the committed prefix and commits the longest prefix with $y_i=F_i(y_{<i})$ plus the target's token at the first mismatch. The commit is exact for any candidate, because every committed token is the target's own decision on a verified prefix and the committed recurrent state and attention cache are those the target computed for that prefix. The Jacobi map $y^{(r+1)}_i=F_i(y^{(r)}_{<i})$ fixes at least one more position per application and reaches the solution within the block length; this is Jacobi iteration applied to decoding~\cite{song2021parallel,santilli2023jacobi}, and lookahead and consistency-trained decoders pursue related fixed-point iterations~\cite{lookahead,cllm}. For sampling, fixed per-position random inputs held across repairs give the target law, as in Theorem~\ref{thm:fixednoise}.

Let DFlash commit $A_D$ tokens per cycle of cost $C_D$ (draft, verification, state commit and control). A repair scheme that commits $A_R$ tokens per attempt at cost $C_R=C_{\rm init}+C_{\rm anchor}+k\,C_{\rm repair}+C_{\rm audit}+C_{\rm state}+C_{\rm control}$, averaged over all attempts including failures and fallbacks, speeds decoding up by $S_{\rm decode}=A_RC_D/(A_DC_R)$. Prefill and scheduling take a fraction $f$ of a request's time that repair does not change, so end to end $S_{\rm e2e}=1/(f+(1-f)/S_{\rm decode})$; a $5\times$ end-to-end gain needs $S_{\rm decode}=5.44$ at $f=2\%$ and 9 at $f=10\%$ (derived). Two oracle bounds use a perfect candidate of width $B$ (the reference continuation, or a forced full acceptance). With $V(B)$ the target verify pass and $C_{\rm state}(B)$ the work needed to continue from the accepted boundary,
\begin{equation}\label{eq:repair-oracle}
S_{\rm a}(B)=\frac{B\,C_D}{A_D\bigl(V(B)+C_{\rm state}(B)\bigr)},\qquad
S_{\rm b}(B)=\frac{B\,C_D}{A_D\bigl(C_{\rm anchor}(B)+C_{\rm audit}(B)+C_{\rm state}(B)\bigr)}.
\end{equation}
$S_{\rm a}$ is an ideal drafter paying one verify pass per block of $B$ tokens; $S_{\rm b}$ is P3's two-pass design with free repair, $C_{\rm anchor}(B)$ a full pass that also writes the anchor cache and $C_{\rm audit}(B)$ a full pass. If $S_{\rm b}$ cannot reach the target end to end, the two-pass design is dead whatever the repair does; if even $S_{\rm a}$ cannot, better drafting alone cannot reach it with this verifier. A concrete repair pays only if $A_R/C_R>A_D/C_D$, and it is rejected outright if a lower bound on its cost gives $C_R^{\rm lower}/(B+1)\ge C_D/A_D$, since no attempt commits more than $B+1$ tokens.

\paragraph{What a repair must achieve}\label{sec:repair-analysis} Local low rank does not imply useful repair of a whole model. Three things must hold separately: the fixed cheap operators stay accurate after candidate tokens change, not only near the anchor; their errors stay below the decision margins after attention, nonlinearities, the recurrent state and later layers; and each sweep resolves enough positions to pay for itself and for the audit. A sufficient condition links these quantities. Let $m_i$ be the exact top-two logit margin at position $i$, $\varepsilon_i$ a bound on the logit error the approximate evaluator makes at $i$ when its inputs are exact, $b_j$ a bound on the error in the current guess's effect at an earlier position $j$ ($b_j=0$ once $j$ is correct), and $C_{ij}$ the influence of position $j$ on the logits at $i$. If $2(\varepsilon_i+\sum_{j<i}C_{ij}b_j)<m_i$, position $i$ is correct after the next sweep, by the triangle inequality. Two consequences follow. A persistent approximation error $\eta$ in a recurrence with contraction factor $a<1$ leaves an error floor of about $\eta/(1-a)$. And a faithful evaluator can advance only as fast as exact Jacobi iteration, which on a copy chain $F_i(y)=y_{i-1}$ gains one position per sweep even with exact rank-one operators; a non-faithful map can jump ahead, as consistency-trained models do~\cite{cllm}, but that would be better drafting, not repair. Expected accepted length is best read through hazards: with $h_i$ the conditional failure rate at position $i$ given an accepted prefix, $\E A=\sum_a\prod_{i\le a}(1-h_i)$. The decisive control is a small target-conditioned drafter given the same anchor information and matched in compute, memory traffic, training data, verifier width and target calls; if it does as well, the supported claim is better target-conditioned drafting, not cheap repair of a trajectory. This is analysis that fixes the measurements, not a result.

\paragraph{Test 1: perfect continuations at width $B$}\label{sec:repair-stagea} $V(B)$ and $C_{\rm state}(B)$ are measured inside SGLang's DFlash worker in the serving configuration (FlashInfer attention and GDN kernels, CUDA graphs, overlap scheduling, concurrency 1) at block widths $B=16,32,64,128,256$, forcing every verify pass to accept the whole block; the committed tokens are then meaningless, but a pass's cost does not depend on token values. CUDA events around each phase of the cycle give GPU times without host synchronization. SGLang's GDN verify kernels write one FP32 state per verified position and the commit copies the accepted one back; the replay protocol, which keeps per-token inputs and folds the accepted prefix at commit, and a microbenchmark that separates the per-position state writes from the verify arithmetic are timed alongside. Feeding the reference continuation instead of forcing acceptance shows how often a verify pass at width $B$ rejects the plain-decoding continuation; such rejections are the stock noise floor of Appendix~\ref{app:divergence}, not drafting errors. $V(B)$, $C_{\rm state}(B)$, the draft cost and cycle period per width, the acceptance of the reference continuation per width, and $S_{\rm a}(B)$ and $S_{\rm b}(B)$ are \pend{repair-a}.

\paragraph{Test 2: Jacobi repair from DFlash windows}\label{sec:repair-jacobi} The second test asks how close exact repair from real DFlash windows gets, with no reference tokens. Every DFlash block $y^{(0)}$ is verified, and then $k$ further full target passes run from the same committed prefix: \emph{recycle} sweeps apply the Jacobi map, and \emph{keep} sweeps replace only the first mismatching token by the target's token and keep the rest of the draft. The pass on the original draft is then re-run and committed, so the trajectory is plain DFlash and each sweep's accepted length is measured on real windows. Re-running a verify pass on the same prefix needs the committed GDN convolution window restored first, because the verify kernel shifts it in place, and the probe checks that the re-run reproduces the first pass's argmax exactly. A sliding variant drafts each new block from the previous pass's suffix predictions and fills the rest with DFlash. Accepted drafts against the number of sweeps, recycle and keep, per block width, are \pend{repair-jacobi}.

\paragraph{Test 3: anchored residual evaluation}\label{sec:repair-residual} For every linear operator $z=Wx$, P3 keeps the anchor $(\bar x,\bar z=W\bar x)$ from a full pass over a DFlash block and repairs a changed block with $\tilde z=\bar z+(WU)U^\top(x-\bar x)$, where the orthonormal basis $U$ is fixed per operator input and $WU$ is precomputed, so a repair sweep streams $WU$ instead of $W$; a basis computed from the current correction would stream $W$ again. Nonlinearities, normalizations, the causal convolution, the GDN recurrence from the exact committed state and attention over the exact prefix run exactly on the repaired values. Bases are the leading principal components of each operator's exact input changes on a development split (split by request), at ranks 16 to 128. Each held-out block is replayed twice against the original anchor, $y^{(1)}=F(y^{(0)})$ and $y^{(2)}=F(y^{(1)})$, and at every position the test records whether $\argmax\tilde z_i=\argmax z_i$, the ratio $R_i=2\norm{\tilde z_i-z_i}_\infty/m_i$ ($R_i<1$ guarantees the same decision), the first disagreement, and the change from the first replay to the second. A free-running repair then iterates, audits every iterate exactly, and reports the accepted prefix after each sweep and the hazards $h_i$. Decision agreement, $R_i$, hazards and progress per rank, and the kill-test verdicts for P2 and P3 against the criteria above (with the first of the three P3 conditions to fail) are \pend{repair-residual,repair-verdict}.\label{sec:repair-verdict}

\subsection{Selector and depth analysis not tested here}\label{app:untested}
Three pieces of analysis bear on the proposals above; none is under test.

\paragraph{Executing the DFlash~2 selector} At the inspected implementation adjacent candidate scores have the form
\begin{equation}
S_t(a,b)=U_t(b)+\sum_{r=1}^{d_s}A(a)_r\,H(h_t)_r\,B(b)_r,\label{eq:selector}
\end{equation}
with a gate that depends on the draft hidden state~\cite{dflash,sgcode}. A greedy walk needs only the realized row $S_t(x_{t-1},\cdot)$, which reduces score work from $O(BLK^2d_s)$ to $O(BLKd_s)$ at the price of a dependent sequence of small contractions. Alternatively the lattice can emit, for every predecessor, its winning successor $f_t(a)=\argmax_bS_t(a,b)$, after which the path is $x_t=(f_t\circ\cdots\circ f_1)(x_0)$.
\begin{proposition}[Finite-map execution]\label{prop:maps}
Given the same row winners, serial map walking and prefix composition produce the same path. A work-efficient scan of length $L$ over $K$-entry maps uses $O(LK)$ map work and $O(\log L)$ composition stages, excluding the score producer.
\end{proposition}
\begin{proof}
The $t$-th prefix composition applies exactly the first $t$ transitions to the anchor, and function composition is associative. A work-efficient tree scan uses $O(L)$ compositions of $O(K)$ indexed evaluations each, with logarithmic depth.
\end{proof}
Function scanning is standard; the open question would be where it pays. The exact rational reference implements the simpler $O(LK\log L)$ Hillis--Steele scan and checks path equivalence only (Table~\ref{tab:tests}).

\paragraph{Prefix value on a fixed lattice} With a substochastic surrogate $r_t(a,b)$ for the probability of matching the next target token, $J_r(x)=\sum_{j=1}^{L}\prod_{t\le j}r_t(x_{t-1},x_t)$ is maximized by the recursion $V_t(a)=\max_br_t(a,b)(1+V_{t+1}(b))$ with $V_{L+1}=0$. If the estimate $\hat r$ also lies in $[0,1]$ and edge errors are uniformly bounded by $\epsilon_t$, telescoping gives $|J_r-J_{\hat r}|\le\delta=\sum_t(L-t+1)\epsilon_t$ and optimizer regret at most $2\delta$. LiLiCorr's finding that an exact best-path search over a different score decodes worse than the greedy rule~\cite{lilicorr} is the relevant warning.

\paragraph{Choosing depth under a cost model} If predicted progress $F_i(d)$ is known and cost depends only on the total optional depth $m$, the depth assignment for a batch is solved by
\begin{equation}
G_i(m)=\max_{0\le d\le L_i}\{G_{i-1}(m-d)+F_i(d)\},\qquad\max_m\frac{G_B(m)}{C(m;z)},\label{eq:schedDP}
\end{equation}
which is exact only under that cost hypothesis. Published policies choose depth in several ways: OPT-Tree searches for the tree that maximizes expected acceptance~\cite{opttree}, Sequoia sizes trees for the hardware~\cite{sequoia}, SpecDec++ stops drafting with a learned acceptance head~\cite{specdecpp}, DSpark schedules verification from confidences and a profiled throughput curve~\cite{dspark}, D-cut prunes draft blocks against a profiled cost table aligned with the engine's CUDA-graph shapes~\cite{dcut}, and an open SGLang pull request ports the DSpark scheduler to DFlash~2 with greedy benchmarks~\cite{pr36136}. CUDA-graph tiers make the cost discontinuous: with rewards $(1,1.9,2.7)$ and costs $(1,2,2)$ at depths $(0,1,2)$, depth one looks worse than depth zero while depth two is best, so stopping at the first local decline fails. Autotuning a kernel that mutates state must restore that state between trials~\cite{tritonautotune}. Any such policy must choose depth from committed-prefix information when sampling (Appendix~\ref{app:stopping}).
